\section{PyMARL}
\label{sec:pymarl}
To make it easier to develop algorithms for SMAC, we have also open-sourced our software engineering framework PyMARL. PyMARL has been designed explicitly with deep MARL in mind and allows for out-of-the-box experimentation and development.

PyMARL's codebase is organized in a modular fashion in order to enable the rapid development of new algorithms, as well as providing implementations of current deep MARL algorithms to benchmark against. It is built on top of PyTorch to facilitate the fast execution and training of deep neural networks, and take advantage of the rich ecosystem built around it. PyMARL's modularity makes it easy to extend, and components can be readily isolated for testing purposes.

Since the implementation and development of deep MARL algorithms come with a number of additional challenges beyond those posed by single-agent deep RL, it is crucial to have simple and understandable code. In order to improve the readability of code and simplify the handling of data between components, PyMARL encapsulates all data stored in the buffer within an easy to use data structure. This encapsulation provides a cleaner interface for the necessary handling of data in deep MARL algorithms, whilst not obstructing the manipulation of the underlying PyTorch Tensors. In addition, PyMARL aims to maximise the batching of data when performing inference or learning so as to provide significant speed-ups over more naive implementations. 

PyMARL features implementations of the following algorithms: QMIX \cite{rashid_qmix:_2018}, QTRAN \cite{son2019qtran} and COMA \cite{foerster_counterfactual_2017} as state-of-the-art methods, and VDN \cite{sunehag_value-decomposition_2017} and IQL \cite{tan_multi-agent_1993} as baselines.




